% Abstract
This dissertation presents \textbf{MutBench}, a hotspot-detection benchmark that systematically compares 10 biological information types (20 scoring formulas) $\times$ 14 detection families (39 variants) $\times$ 11 RNA viruses = 8{,}580 evaluations. The headline result: the optimal information source is pathogen-dependent, with the scoring~$\times$~pathogen interaction ($\omega^2 = 0.296$, 11-pathogen cluster-bootstrap 95\% interval [0.195, 0.333]) as the largest modeled variance component, and HIV-1 vaccine-escape enrichment of $7.19\times$ (Bonferroni $p = 2.5 \times 10^{-16}$) as the primary external anchor.

Detecting mutation hotspots in viral genomes is essential for variant surveillance and vaccine design, yet existing approaches have relied almost exclusively on mutation frequency and Shannon entropy, leaving comparatively unexplored a rich landscape of complementary information sources (phylogenetic selection signals, protein language models, structural features). Concurrent fitness-regression benchmarks (ViroGym, EVEREST~v3, ProteinGym~v1.3) target per-variant prediction; MutBench is the broadest \emph{region-level hotspot-detection} benchmark to date.
MutBench comprises three core components:
(1) a three-layer ground truth system integrating literature-curated positive markers (Layer~A: a pathogen-specific union of convergent-evolution, immune-escape, and HVR-membership evidence types, with immune-escape positions dominating the curated set), conserved structural regions under purifying selection as negative controls (Layer~B), and Deep Mutational Scanning (DMS) experimental fitness data as functional validation (Layer~C);
(2) a composite evaluation metric defined as $\text{hotspot-score} = (\text{recall} + \text{precision} + \text{stability}) / 3$; and
(3) a two-stage experimental design in which Stage~1 conducts in-depth analysis on SARS-CoV-2 and Stage~2 performs the large-scale cross-pathogen comparison.

Stage~1 validates the framework on SARS-CoV-2 (MutClust-Hybrid hotspot-score 0.778). Stage~2 establishes that the optimal information source varies systematically across pathogens: $\omega^2_{\text{scoring}\times\text{pathogen}} = 0.296$ (11-pathogen cluster-bootstrap 95\% interval [0.195, 0.333]; 6-category aggregation gives $\omega^2 = 0.103$ as a collinearity-robust lower bound; within-cell residual $\omega^2 \approx 0.39$). Friedman test shows no universally best combination ($\chi^2 = 7.69$, $p = 0.990$), and LOPO yields 0/11 exact matches with a 0.265 oracle-vs-generalized MCC gap (LOPO 0/11 alone is null-consistent under permutation; the conclusion rests on the interaction effect plus the vaccine-escape audit, not on LOPO in isolation).
Contribution 3 is two-part: \textit{(3a) Multi-source integration infrastructure} ---a 4-feature core captures $\approx$98\% (97.6\%) of the 10-feature ensemble; EqualWeight integration achieves H3N2 enrichment 4.012$\times$ ($p = 0.0023$). \textit{(3b) External validation, anchored on HIV-1} ---under a three-tier evidence rule (Bonferroni survival~$\times$~Layer~A overlap < 33\%), HIV-1 is the primary anchor with direct novel-only Fisher enrichment $7.16\text{--}8.24\times$ on the 37 Layer-A-disjoint positions (worst-case $p = 5 \times 10^{-12}$); H3N2 ($9.36\times$, 69\% overlap) is a self-consistency check; SARS-CoV-2 ($7.63\times$) is exploratory after Bonferroni. Validation is restricted to 3/11 pathogens (2{,}340 evaluations) by availability of curated antibody-escape annotations.

A systematic analysis of the 20 scoring types spanning six information categories (frequency, MSA-derived, phylogenetic, structural, AI-based, and composite) reveals 9 distinct optimal information types across 11 pathogens: FUBAR phylogenetic selection scores achieve the highest MCC for H3N2 (0.534), protein language model scores dominate for SARS-CoV-2 and RSV, and homoplasy scores excel for Norovirus.
Multi-source integration further reduces the experimental search space for hotspot validation, achieving 4.012$\times$ vaccine escape enrichment for H3N2 ($p = 0.0023$), while a compact 4-feature core captures $\approx$98\% (97.6\%) of full-ensemble performance.

\vspace{1em}
\noindent\textbf{Keywords}: mutation hotspot detection, MutBench, biological information types, viral genomics, multi-source information analysis, pathogen-dependent optimality, SARS-CoV-2, Deep Mutational Scanning
